\section{Uncertainty Propagation} \label{sec:uncertainty_propagation}

As introduced in \cite{QA4EO2}, there are two methodologies to propagate
uncertainties through the measurement model until the measurand. The first
method is the \ac{LPU}, which relies on the error covariance matrices and
sensitivity coefficient produced for each output quantity. The second method is
the \ac{MC} method, which relies on the probability distribution function of
each source of uncertainty or, in a compact way, the output quantities. This
section details how to apply these methods during the uncertainty analysis.

\subsection{Law of propagation of uncertainties method}

The \ac{LPU} method is based on a locally linear
approximation to combine uncertainties, expressed as

\begin{equation}
  u^2(y) = \sum_{i=1}^N c_i^2 u^2(x_i) + 2 \sum_{i=1}^{N-1} \sum_{j=i+1}^N c_i c_j u(x_i, x_j).
  \label{eq:propagation}
\end{equation}

\noindent Here, $u^2(y)$ is the combined uncertainty associated with the
measured value $y$. $u(x_i)$ is the uncertainty associated with the input
quantity $x_i$. $u(x_i, x_j)$ is the covariance between $x_i$ and $x_j$.
$c_i=\delta y/ \delta x_i $ is the sensitivity coefficient, i.e., the
relationship between an uncertainty associated to an input quantity and the
uncertainty associated to the measured quantity. $N$ represents the number of
input quantities. In this equation the first term, $\sum_{i=1}^N c_i^2
u^2(x_i)$ is fully independent, while the second term $2 \sum_{i=1}^{N-1}
\sum_{j=i+1}^N c_i c_j u(x_i, x_j)$ stands for the error correlation.

The expression can also be written in matrix form as $u^2(y)=cS(x)c^T$, where
$c=[\delta y/\delta x_1, ..., \delta y/\delta x_N]$ is a row vector of
sensitivity coefficients and $S(x)$ is the error covariance matrix for the
input quantities with diagonals $u^2(x_i)$ and off-diagonals $u(x_i, x_j)$.

\subsection{Monte Carlo method}

The \ac{MC} method estimates the input probability distributions by finite sets
of random draws from the those distributions. These draws are processed
through a measurement function to produce output values, which are treated as
random samples of the output probability distribution. Statistical analysis,
such as calculating standard deviations and error covariances, is then
performed on the output values. 

The \ac{MC} method offers advantages such as providing insights into the shape
of the output probability distribution, handling highly non-linear measurement
functions, and accommodating complex probability distributions. They are
particularly useful for models that cannot be expressed algebraically. However,
they are computationally intensive, which is a significant drawback when
dealing with large Earth Observation data volumes. In practice, uncertainty
analysis often combines \ac{MC} with \ac{LPU}, where \ac{MC} is used to
determine uncertainty for specific quantities, which are then utilized in
\ac{LPU} for subsequent analysis.
